\section{Pure critical Ising inputs and contour conventions}
\label{sec:pure}

The comparison and regularity estimates in this section concern only the
nearest-neighbor Ising model at coupling
\(K_0=\beta_0J=\frac12\log(1+\sqrt2)\). Their role is to supply uniform
control of bounded spin observables across a free region, and geometric
estimates for the contours of the pure model. Section~\ref{sec:transfer}
will derive the uniform transfer theorem from these inputs.

We write \(\mu_0\) for the critical full-plane law and \(\E_0\) for its
expectation. A \emph{pin} is a prescribed spin. Common pins have the same
locations and values in the laws being compared; they may have either
sign. A free lattice chart is a region isometric to part of the square
lattice in which no spins have been pinned. A collar always includes
the edges of that chart, as well as its vertices. Write
\(B_s(x)=x+[-s,s]^2\cap\Z^2\), and \(B_s=B_s(0)\). For integer \(s\),
the outer layer is \(\partial B_s=B_s\setminus B_{s-1}\); integrating
the open square means integrating \(B_{s-1}\) while keeping this layer
fixed. Integer rounding and a bounded displacement of these layers can
be absorbed by a fixed fraction of any stated positive collar once the
lattice scale tends to infinity.

We cite three equilibrium estimates and one complete-resolution result
from the companion article \emph{Buffered comparison and stopping-band
resolution in critical Ising}~\cite{OpenAIStatic2026}. The citations
below specify its named results and section headings. We use these
results as pure-model inputs; the
transfer, scale flow, and perturbed-interface arguments are proved in
this article.

\subsection{Pointwise comparisons through collars}

Pointwise comparison is stronger than comparison of a fixed collection
of observables. It will permit arbitrary boundary assignments to enter
the transfer argument through uniformly bounded likelihoods.

\begin{proposition}[Buffered and small-patch comparisons]
\label{prop:pure-buffered}
Consider a finite critical nearest-neighbor ferromagnetic Ising graph,
disjoint vertex sets \(I,J,F\), and common pins on \(F\). Let \(\mu_b\)
be the marginal law on \(I\) obtained by prescribing an arbitrary
configuration \(b\) on \(J\). Whenever the following sets lie in the
indicated square-lattice charts, the bounds hold for all configurations
\(a\) on \(I\) and all \(b,b'\).
\begin{enumerate}
\item For every \(c>0\) there is \(K_c<\infty\) such that
\(I\subset B_R(z)\), \(J\subset B_{(1+c)R}(z)^c\) imply
\[
 K_c^{-1}\le \frac{\mu_b(a)}{\mu_{b'}(a)}\le K_c.
\]
\item There are absolute \(K<\infty\), \(\zeta>0\) such that
\(I\subset B_w(z)\), \(J\subset B_R(z)^c\), and \(R\ge8w\ge8\)
imply
\[
 \left|\log\frac{\mu_b(a)}{\mu_{b'}(a)}\right|
 \le K(w/R)^\zeta.
\]
No free neighborhood of \(I\) is required in these two assertions:
additional common pins may lie anywhere outside \(I\cup J\).
\item Either law can be replaced by any mixture of the indicated
conditional laws. The fixed-ratio comparison also holds for laws from
different domains, including the plane and a torus, if there is a common
isometric lattice collar separating the observed spins from every
difference in the data. In this cross-domain assertion the collar is
unpinned in both laws. Applying a common randomization kernel to the
observed spins preserves each density bound.
\end{enumerate}
All constants are independent of the number, positions, and values of
the common pins and of the mesh. In the first two assertions the same
ratio bounds hold when the roles of the two terminal sets are reversed.
\end{proposition}

\begin{proof}[Source and the terminal symmetry]
The first three assertions are Corollary~2.5, ``Buffered and small-patch
comparisons'', in Section~2, ``Pointwise comparison through buffers'',
of~\cite{OpenAIStatic2026}. We record the symmetry needed below. For the
positive joint marginal \(p(a,b)\), put \(f=\log p\). The source's
Lemmas~2.1 and~2.2, ``Extreme rectangles in a ferromagnetic marginal''
and ``Pointwise ratios with arbitrary common pins'', bound
\[
 \left|f(a,b)+f(a',b')-f(a,b')-f(a',b)\right|
 \le 4\operatorname{arctanh}q_0,
\]
where \(q_0\) is the zero-field FK connection probability after deleting
the common pinned vertices and separately contracting the two terminal
sets. Deletion retains the fields induced by the pinned spins in the
original law; all these fields are then set to zero in the comparator.
The right side is symmetric in the terminals. A dual circuit in
the separating collar bounds \(q_0\) away from one; several separated
dual circuits give its power decay in the small-patch case. Normalizing
either family of conditional laws turns the oscillation bound into the
displayed pointwise ratio bounds. Thus reversing which terminal is
observed requires no new geometric estimate; this is also
Remark~2.3 of the companion.
\end{proof}

\begin{corollary}[Two collars around an annulus]
\label{cor:pure-annular}
Let \(2\le r<R/4\),
\(A=B_R\setminus B_r\), and \(W=B_{2R}\setminus B_{r/2}\), with
integer rounding of the boundaries. Suppose \(W\) is an unpinned
square-lattice chart. The law on \(A\), conditional on arbitrary spins
outside \(W\), has pointwise density relative to \(\mu_0|_A\) between
\(K^{-1}\) and \(K\), for an absolute \(K\) independent of \(R/r\).
The assertion applies in a torus chart and to mixtures of exterior
conditions. Fixed changes of the two collar ratios only change \(K\).
\end{corollary}

\begin{proof}
First change the spins outside \(B_{2R}\), keeping all inner pins
common. Proposition~\ref{prop:pure-buffered} across the outer collar
costs a fixed factor. Next change the spins in \(B_{r/2}\), keeping the
outer data common. Use the terminal-symmetric form of the same
proposition across the collar from \(r/2\) to \(r\): the small terminal
is now the source and \(A\) lies beyond its fixed-ratio collar. This
costs another fixed factor. The product is independent of the distance
between the collars. Finally mix the exterior data according to the
full-plane law. Cross-domain comparison is legitimate because both
collars are common lattice charts.
\end{proof}

\subsection{Signed sources and line magnetizations}

The transfer proof will produce a signed density that has a positivity
property under increasing tests, but is not necessarily increasing.
The following definition states precisely the property required of it.

\begin{definition}[Order-positive source]
\label{def:pure-order-positive}
For a probability measure \(\mu_S\) on spins in a finite set \(S\), a
real function \(q\) is an order-positive centered density if
\[
 \E_{\mu_S}q=0,\qquad \E_{\mu_S}[qg]\ge0
 \quad\hbox{for every increasing }g.
\]
Here density means the density of a signed measure.
\end{definition}

We must also specify the small free neighborhoods used by the
transmission estimate. Let \(S,T\) be disjoint source and target sets.
A cut \(J\), disjoint from \(S,T\) and from all common pins, is
\emph{separating} if every free-graph path from \(S\) to \(T\) meets
\(J\). Its target-side Gibbs continuation is then independent of the
source after the cut spins have been prescribed.

Fix \(A,K\ge1\) and \(a\in(0,1/(16A))\). The cut is
\emph{admissible at scales \(w,s\)}, where \(0<w\le s\), if it has a
partition into patches \(b_i\), straight counting segments
\(L_i\subset J\), centers \(z_i\), and scales
\(\ell_i\in[aw,Aw]\) such that
\begin{enumerate}
\item \(b_i\cup L_i\subset B_{\ell_i}(z_i)\) and
\(a\ell_i\le |L_i|\le2\ell_i\);
\item \(B_{16\ell_i}(z_i)\) is a pin-free planar lattice chart,
including freedom from source pins;
\item each vertex belongs to at most \(K\) counting segments;
\item \(\dist(T,z_i)\ge s/A\), \(16\ell_i\le s/A\), and the
small-patch comparison applies from \(b_i\) toward \(T\) in a chart
extending to a fixed positive multiple of \(s\).
\end{enumerate}
The constants \(a,A,K\) and the fixed chart geometry are part of this
definition; the number of patches is not. A square contour in a free
collar of width comparable to \(w\), at distance comparable to \(s\)
from the target, is admissible. Partition its sides into intervals of
length a sufficiently small fixed multiple of \(w\). At corners use
counting intervals on an adjacent side. Their overlap remains bounded.
The same construction works for a fixed finite union of square contours,
including the two boundaries of a square band.

\begin{proposition}[Transmission of a signed source]
\label{prop:signed-transmission}
There is a fixed exponent \(\eta\in(0,1/8)\) with the following
property. Let \(\mu\) be a pure critical Gibbs law with common pins,
let \(q=q(\sigma_S)\) be order-positive and centered for its source
marginal, and let \(J\) be an admissible cut at scales \(w,s\).
For every bounded target observable \(h=h(\sigma_T)\), when \(w\) is
sufficiently large,
\begin{equation}
 \left|\E_\mu[qh]\right|
 \le C\norm{h}_\infty (w/s)^{-7/8+\eta}
               \E_\mu[qM_J],
 \qquad M_J=s^{-7/8}\sum_{x\in J}\sigma_x.
 \label{eq:pure-signed}
\end{equation}
The constant depends only on the admissibility constants and the fixed
chart geometry. It is independent of the mass of \(q_+\), the mesh,
the number and values of source pins, and the ratio \(w/s\).
The same inequality holds for expectation differences of two ordered
source configurations, and after averaging an ordered coupling of
random source configurations, with the common Gibbs continuation.
All changed pins and fields lie on the source side of the cut; the
target-side continuation is identical in the compared laws.
One may decrease \(\eta\), keeping it positive.
\end{proposition}

\begin{proof}[Source and signed-measure interpretation]
These are Proposition~4.2, ``Transmission'', and Corollary~4.11,
``Random and signed sources'', in Section~4, ``Ordered transmission
across a cut'',
of~\cite{OpenAIStatic2026}. To clarify the signed assertion,
put \(m=\E q_+=\E q_-\). If \(m=0\), both sides vanish. Otherwise
the source probability measures \(m^{-1}q_+\mu_S\) and
\(m^{-1}q_-\mu_S\) are stochastically ordered by
Definition~\ref{def:pure-order-positive}. Couple them in that order,
apply the deterministic-source transmission inequality, and multiply
by \(m\). Its nonnegative magnetic difference permits averaging and
gives~\eqref{eq:pure-signed}. No division by a mesh-dependent lower
bound on \(m\) occurs.
\end{proof}

\begin{proposition}[Nearby line sums]
\label{prop:line-continuity}
Let \(J_s,J'_s\) be horizontal or vertical lattice segments of length
\(O(s)\), with corresponding portions displaced by at most \(Cw\)
and unmatched end portions of total length \(O(w)\). The assertion
also applies to concentric square contours of radii comparable to
\(s\), differing in radius by at most \(Cw\), and to a fixed finite
union of these pairs. Assume their union has a free neighborhood of
width \(cs\), covered by a fixed number of buffered lattice charts.
For each fixed \(h=w/s\in(0,h_0]\),
\begin{equation}
 \limsup_{s\to\infty}
 \left\|s^{-7/8}\sum_{J_s}\sigma_x
       -s^{-7/8}\sum_{J'_s}\sigma_x\right\|_{L^2(\mu)}
 \le C_0h^{3/8}.
 \label{eq:pure-lines}
\end{equation}
The bound is uniform over arbitrary pins beyond the stated free
neighborhood and mixtures of their laws. The constant depends only on
the fixed geometry and buffer.
\end{proposition}

This is Proposition~5.1, ``Differences of nearby line sums'', in Section~5,
``Nearby contours and prediction'', of~\cite{OpenAIStatic2026}. Its exponent follows from integrating
the full-plane covariance kernel \(c|x-y|^{-1/4}\) along segments:
separations below \(h\) contribute \(O(h^{3/4})\) to the squared
norm, while above \(h\) the first-derivative bound gives
\(Ch\int_h^1t^{-5/4}\,dt=O(h^{3/4})\). Short unmatched end pieces
contribute \(O(h^{7/4})\) in squared norm. Buffered pointwise
comparison transfers the nonnegative squared error to the allowed
laws. This explanation uses the continuum kernel after taking the
mesh limit and asserts no uniform discrete derivative bound.

\begin{corollary}[Prediction from a nearby contour]
\label{cor:line-prediction}
In Proposition~\ref{prop:line-continuity}, let \(\mathcal F\) be any
spin information containing \(\sigma_{J'_s}\), and let
\(P_{\mathcal F}\) denote conditional expectation. Then, for each
fixed \(w/s\),
\[
 \limsup_{s\to\infty}
 \norm{M_{J_s}-P_{\mathcal F}M_{J_s}}_2
 \le C_0(w/s)^{3/8}.
\]
Both line sums use the same normalization \(s^{-7/8}\).
\end{corollary}

\begin{proof}
The second line sum is an \(\mathcal F\)-measurable competitor for the
orthogonal projection. Apply~\eqref{eq:pure-lines}. This is also the
Corollary~5.2, ``Prediction from a nearby contour'', in Section~5 of
\cite{OpenAIStatic2026}.
\end{proof}

In particular, when \(w\asymp s/n\) for \(n\) concentric cuts,
Proposition~\ref{prop:signed-transmission} has a constant independent
of \(n\). In Proposition~\ref{prop:line-continuity} one first fixes
\(n\) and then takes \(s\) large; the threshold may depend on
\(n\). The fixed exponent \(\eta>0\), geometric constants, and
normalization are chosen before either operation.

\subsection{Spin correlations and the thermal direction}

\begin{proposition}[Spin normalization and moment bounds]
\label{prop:pure-spin}
There is \(a_\sigma>0\), determined by the square-lattice convention,
such that, for distinct physical points \(z_1,\ldots,z_n\),
\[
 a_\sigma^{-n}\delta^{-n/8}
 \E_0\prod_{i=1}^n\sigma_{x_i^\delta}
 \longrightarrow S(z_1,\ldots,z_n),
 \qquad \delta x_i^\delta\longrightarrow z_i.
\]
For odd \(n\), \(S=0\); for even \(n\), \(S\) is the nonnegative
square root of
\begin{equation}
 S(z_1,\ldots,z_n)^2
 =2^{-n/2}\!\sum_{\substack{u_i\in\{-1,1\}\\\sum_i u_i=0}}
       \prod_{i<j}|z_i-z_j|^{u_iu_j/2}.
 \label{eq:pure-spin-formula}
\end{equation}
Thus \(S(z,w)=|z-w|^{-1/4}\). The convergence is uniform when the
points lie in compact sets with fixed mutual separation. Also
\begin{align}
 0\le\E_0[\sigma_x\sigma_y]&\le C(1+|x-y|)^{-1/4},
 \label{eq:pure-two-point}\\
 0\le\E^+_{B_r}[\sigma_0]&\le C(1+r)^{-1/8}.
 \label{eq:pure-plus-one-point}
\end{align}
For \(n\ge2\), repeated sites allowed,
\begin{equation}
 \left|\E_0\prod_{i=1}^n\sigma_{x_i}\right|
 \le C^n\prod_{i=1}^n
 \left(1+\min_{j\ne i}|x_i-x_j|\right)^{-1/8}.
 \label{eq:pure-multipoint-bound}
\end{equation}
For every fixed finite collection of bounded compactly supported
Riemann-integrable spatial test functions \(\varphi\), the smeared
variables
\[
 \Phi_\delta(\varphi)
 =a_\sigma^{-1}\delta^{15/8}
                  \sum_{x\in\Z^2}\varphi(\delta x)\sigma_x
\]
have joint moment limits obtained by integrating \(S\) against the
test functions. Here Riemann integrability is equivalent to having a
set of discontinuities of two-dimensional Lebesgue measure zero.
Spatial smearings are chosen in this class and fixed before taking
the mesh limit; polygonal cell indicators are included.
The limiting correlation functions are covariant under
M\"obius transformations, including inversion, with weight \(1/8\)
at each spin insertion.
\end{proposition}

\begin{proof}[Sources and the domination needed for smearing]
The construction of smeared magnetization fields from spin correlations
is developed in \cite[Section~3]{CGN2015}.  We give the domination
argument for the spatial test functions used here.
The spin limit and its full-plane normalization follow from
\cite[Theorems~1.1--1.2, Remark~1.4, and Equation~(1.6)]{CHI2015}.
That paper uses a rotated lattice with
nearest-neighbor spacing \(\sqrt2\) times its mesh parameter; its
positive normalization constant is absorbed into \(a_\sigma\).
The plus-boundary one-point convergence in a fixed square and the
full-plane two-point asymptotic give~\eqref{eq:pure-plus-one-point}
and~\eqref{eq:pure-two-point} for large distances. Enlarging \(C\)
covers the finitely many smaller distances.

For~\eqref{eq:pure-multipoint-bound}, put disjoint squares of radius
a small fixed fraction of the nearest other insertion distance around
the distinct insertion sites. Conditional on the spins outside these
squares, the squares are independent. Boundary monotonicity bounds the
absolute conditional mean at each insertion by the corresponding
plus-square mean. A site appearing an even number of times contributes
one; for a repeated site the corresponding factors on the right of
\eqref{eq:pure-multipoint-bound} are one. Multiplication gives the
bound, after harmless bounded-distance changes of the boxes.

For each fixed \(n\ge2\), the rescaled correlation kernel
\[
 K_{\delta,n}(z_1,\ldots,z_n)
 =a_\sigma^{-n}\delta^{-n/8}
          \E_0\prod_{i=1}^n\sigma_{z_i/\delta},
 \qquad z_i\in\delta\Z^2,
\]
is uniformly bounded in \(L^2\) on bounded products, with normalized
lattice counting measure \(\delta^{2n}\). To see this, square the
nearest-distance majorant in~\eqref{eq:pure-multipoint-bound} and
bound the product of nearest-neighbor maxima by
\[
 |K_{\delta,n}(z_1,\ldots,z_n)|^2
 \le C_n\sum_{\substack{f:\{1,\ldots,n\}\to\{1,\ldots,n\}\\
                       f(i)\ne i\text{ for all }i}}
       \prod_{i=1}^n(\delta+|z_i-z_{f(i)}|)^{-1/4}.
\]
There are finitely many directed choices \(f\). Each resulting graph
consists of directed cycles with attached trees. Integrate tree leaves
first. A two-cycle has a single singularity of power \(1/2\); on a
longer cycle integrate one vertex by Cauchy--Schwarz and then the
remaining path. Only one-variable singularities of power at most
\(1/2<2\) occur. Their bounded-domain integrals and mesh-regularized
lattice sums are uniformly bounded, including when lattice insertion
sites coincide. This proves the \(L^2\) bound; its constants may
depend on the fixed \(n\), with no assertion here about their growth
in \(n\). The same integration argument with \(\delta=0\) gives
the corresponding bound for the continuum correlation kernel. The
first moment vanishes by spin symmetry, and the zeroth moment is one.

On a fixed bounded product, the set where any two insertion points
are within distance \(\varepsilon\) has normalized counting measure
\(O_n(\varepsilon^2+\delta^2)\). Cauchy--Schwarz therefore bounds
its contribution to a smeared moment by
\(C_n(\varepsilon+\delta)\), with the fixed test-function bounds
absorbed in \(C_n\). Away from these diagonals the correlation
convergence is uniform. Multiply by a continuous cutoff that vanishes
when a pair is within distance \(\varepsilon\) and equals one when
all pairwise distances are at least \(2\varepsilon\).
The remaining continuum integrand is bounded and
Riemann integrable: its discontinuities lie in the finite union of
the coordinate lifts of the spatial test functions' null discontinuity
sets. Its lattice Riemann sums consequently converge. First take the
mesh limit and then let \(\varepsilon\downarrow0\), using the
preceding uniform bound. This proves the asserted joint smeared
moments.
Finally the covariance rule follows directly from
\eqref{eq:pure-spin-formula} and
\(|\psi(z)-\psi(w)|=|z-w|\,|\psi'(z)\psi'(w)|^{1/2}\)
for M\"obius maps, with complex conjugation included for an
orientation-reversing map.
\end{proof}

\begin{lemma}[Reflection positivity]
\label{lem:pure-reflection}
The full-plane law \(\mu_0\) is reflection positive across a horizontal
or vertical lattice reflection line. The limiting smeared spin
correlation algebra is reflection positive across a straight line.
\end{lemma}

\begin{proof}
The finite-bond argument below is a special case of the
reflection-positivity method of \citet{FILS1978}. In a finite
reflection-symmetric box, reflection through a bond line
leaves the within-half weights paired. Each bond crossing the line has
matrix \((e^{K_0st})_{s,t=\pm1}\), whose eigenvalues are
\(2\cosh K_0\) and \(2\sinh K_0\), both positive. Tensoring its
positive decomposition and summing the spins in the two halves writes
\(\E[F\,\overline{F\circ\vartheta}]\) as a sum of nonnegative
squares for every half-box function \(F\). For a reflection through
vertices, condition on the spins on the line: the two remaining halves
are independent with reflected laws and the same conclusion follows.
Exhaust the plane by symmetric boxes, then use the moment convergence
of Proposition~\ref{prop:pure-spin} for polynomials in smeared fields
whose supports lie strictly on one side. Approximation up to the line,
when needed, follows in \(L^2\) from~\eqref{eq:pure-two-point}.
\end{proof}

\begin{proposition}[Positive energy covariance]
\label{prop:energy}
Let \(e_1,e_2\) be the coordinate unit vectors and set
\(b_j(x)=\sigma_x\sigma_{x+e_j}-1/\sqrt2\). In the full-plane
critical law,
\begin{equation}
 \Cov_0(b_j(0),b_k(x))
   =\frac{1+o(1)}{\pi^2|x|^2},\qquad |x|\longrightarrow\infty,
 \quad j,k\in\{1,2\}.
 \label{eq:pure-energy}
\end{equation}
For a fixed finite local average
\(B(x)=\sum_qc_qb_{j_q}(x+d_q)\), with \(c_q\ge0\) and
\(c=\sum_qc_q>0\),
\[
 \Cov_0(B(0),b_k(x))\sim\frac{c}{\pi^2|x|^2},\qquad
 \Cov_0(B(0),B(x))\sim\frac{c^2}{\pi^2|x|^2}.
\]
In particular square symmetrization of the bond energy cannot cancel
its leading covariance.
\end{proposition}

\begin{proof}[Normalization and exact criticality]
Use only the pure specialization of~\cite[Equations~(1.8), (1.11),
and Theorem~1.2]{GGM2012}. Their scaled energy is
\(a^{-1}\sigma_x\sigma_{x+ae_j}\), and their superscript \(T\)
denotes connected expectation. At two insertions their massless
formula is \(1/(\pi^2|z-w|^2)\), independently of both orientation
labels: the two permutations and the two nonzero sign choices each
give \(|z-w|^{-2}\), with prefactor \(1/(4\pi^2)\).
The remainder bound in their Equation~(1.16), together with the
following propagator estimate, is uniform when the two points remain
in a fixed annulus of separations. This supplies uniformity in the
direction of the displacement.

Their formulation permits a nonzero mass parameter tending to zero.
To obtain the exactly critical lattice limit, fix a sequence \(a\to0\)
and choose a subcritical \(\beta(a)\) so close to \(\beta_0\) that
the finitely many spin covariances being tested differ from their
critical values by \(o(a^2)\), and
\(|\beta(a)-\beta_0|=o(a)\). Such a choice is possible by continuity
of finite-spin expectations at the unique critical Gibbs state.
The cited formula applies along this sequence; the error divided by
\(a^2\) vanishes. This proves~\eqref{eq:pure-energy}. Subtraction of
the mean does not change a covariance. The finite-sum assertions then
follow term by term, since each fixed displacement \(d_q\) is
negligible compared with \(|x|\) and every leading coefficient has
the same positive sign.
\end{proof}

\subsection{Pure interfaces, complete resolutions, and four passages}

A \emph{strong} spin path uses nearest-neighbor edges. A \emph{weak}
spin path may also join vertices of the same lattice face. These
notions must be distinguished at an alternating four-spin face. The
fixed triangulation obtained by joining its northwest and southeast
spin vertices realizes the NE/SW pairing when all disagreement edges
of the face are included. In a bulk square-lattice chart, affine
interpolation on these triangles therefore gives a noncrossing zero
set whose local connections match the rounded dual contours within
\(O(\delta)\), with the same traversal order. The added diagonals
define connectivity and drawing conventions; they do not change the
Ising interaction. At the target boundary, the prescribed open curve
uses only the interior edges and its stubs. Its local drawings beside
signed boundary strips are treated in
Lemma~\ref{lem:curves-local-barrier}.

\begin{theorem}[Pure Dobrushin interface]
\label{thm:pure-sle}
In ordinary lattice approximations of a fixed polygonal Jordan domain
with two distinct marked boundary points, the critical nearest-neighbor
Ising interface with opposite constant signs on the two marked arcs
converges to chordal \(\SLE_3\). Convergence is in uniform distance
modulo increasing reparametrization, oriented from the first mark to
the second. Either local turn at an alternating face is allowed,
including the deterministic rule of the problem. Bounded local
rounding and half-stubs have vanishing effect on this convergence.
\end{theorem}

This is the specialization of~\cite[Theorem~1 and Equation~(1)]{CDHKS2014}
used here. The discussion of ambiguous faces in that paper explicitly
allows arbitrary local turns. Its theorem has a broader domain scope;
only these fixed polygonal approximations are used as direct inputs.
The limiting curve is continuous, simple, and disjoint from the domain
boundary except at its endpoints, by~\cite[Theorems~5.1, 6.1,
and~7.1]{RS2005} and conformal extension in a Jordan domain.

\begin{theorem}[Every complete loop resolution]
\label{thm:pure-cle}
In a plus Jordan domain with lattice boundaries converging uniformly
modulo parametrization, every complete noncrossing resolution of all
spin-disagreement edges converges to the same oriented nested
\(\CLE_3\). The assertion is uniform over choices of the complete
resolution. Loops are compared in uniform distance modulo increasing
reparametrization and choice of starting point; collections are
compared by matching loops and charging the diameters of unmatched
loops. Thus the statement includes coverage and multiplicities of all
loops above each fixed positive diameter cutoff.
\end{theorem}

The published input is~\cite[Theorem~6, its following common-edge
refinement, and Remark~7]{BH2019}. It gives nested convergence for
leftmost and rightmost loops and simultaneous approximation of every
macroscopic oriented edge-simple Ising circuit. The extension to
\emph{every complete resolution} is Lemma~6.1, ``Convergence of every
complete resolution'', in Section~6, ``Complete interface resolutions'',
of~\cite{OpenAIStatic2026}. For clarity, its extra content is not merely a
change of drawing: a pairing switch can merge or split circuits.
The simultaneous circuit approximation and the simplicity and
disjointness of the limiting loops imply that two macroscopic circuits
affected by a switch would have to converge to the same loop. Their
winding numbers would then add to \(0\) or \(\pm2\), contradicting
the winding number \(\pm1\) of the merged simple-loop limit. To fix the
multiplicity, include a specified canonical leftmost circuit in an
auxiliary complete resolution. A second circuit matched to its limit
loop would share an edge with that specified circuit by the common-edge
refinement, contradicting edge-disjointness. The pairing-switch argument
then transfers the count one to every complete resolution. These facts
supply coverage and multiplicity, as required by the stated collection
topology.

A contour \emph{passage} across \(B_R(z)\setminus B_r(z)\) is a
subpath with one endpoint on each of its boundary curves and interior
in the annulus. Several passages use disjoint interiors of parameter
intervals; hit endpoints may be shared. Four passages can belong to
one loop, and repeated counting of one traversal is forbidden. Denote
their existence by \(\mathcal A_4^{\mathrm{cont}}(z;r,R)\).

\begin{proposition}[Iterated four-passage bound]
\label{prop:pure-four-passages}
Fix a bulk lattice chart containing \(B_{2R}(z)\), with all pins and
all changes of ambient domain beyond a fixed proportional outer
collar. For every \(\alpha<21/8\), there are \(C_\alpha<\infty\)
and \(u_\alpha>0\), depending on that fixed collar but not on the
outer conditions, such that, for each fixed \(0<u<u_\alpha\),
\begin{equation}
 \limsup_{\delta\downarrow0}
 \P\bigl(\mathcal A_4^{\mathrm{cont}}(z;uR,R)\bigr)
 \le C_\alpha u^\alpha.
 \label{eq:pure-four-passages}
\end{equation}
This applies to the fixed diagonal resolution, and uniformly to
complete resolutions. The same statement holds for circular annuli
and for fixed-factor changes of the radii. In particular one may fix
\(\alpha>2\). The mesh limit precedes the limit \(u\downarrow0\).
\end{proposition}

\begin{proof}[Conversion from the published spin-arm estimate]
Pointwise comparison reduces the event to a plus-boundary model in a
fixed larger square, with a multiplicative constant depending only
on the outer collar. Rescale so that \(R=1\). For a fixed \(u\),
couple a subsequence of complete loop collections to their limit as
in Theorem~\ref{thm:pure-cle}. Only loops of diameter at least
\(1-u\), or a fixed smaller cutoff after relaxing radii, can contain
the passages under consideration. There are finitely many such limit
loops. Uniform parametrized convergence preserves four passages after
replacing the two radii by \(2u\) and \(1/2\). Indeed choose the
passages by last and first hits of the two circles, pass to convergent
endpoint parameters, and use uniform continuity of the limiting
parametrizations. An interval crossing the relaxed annulus cannot
collapse to a point. Disjoint interval interiors remain distinct;
injectivity of a limiting loop prevents two nontrivial intervals from
having the same image. This also treats passages belonging to one
loop and passages traversed in opposite directions.

We next pass from this limiting geometric event to strong alternating
spin arms. Crop its through subarcs once more, to the radii \(4u\)
and \(1/4\). The finitely many chosen compact subarcs of the simple,
pairwise disjoint limit loops have disjoint thin neighborhoods after
this crop. They are crosscuts of the annulus. Reading their sides in
cyclic order changes the spin sign at each crosscut. If more than four
through crosscuts occur, keep four consecutive ones and the adjacent
alternating sides. One can follow each chosen side in its own thin
neighborhood from the inner to the outer cropped circle.

Now use leftmost approximations for a plus side and rightmost
approximations for a minus side, in the \emph{same} spin configuration.
By~\cite[Section~2.5]{BH2019}, a leftmost loop has a strong plus path
on its left, and a rightmost loop has a strong minus path on its right.
The simultaneous circuit approximation and common-edge refinement
in~\cite[Theorem~6 and Remark~7]{BH2019} make both extremal
approximations converge to the same selected loop, with its orientation.
Their chosen side paths therefore lie in the disjoint neighborhoods
and cross the cropped annulus for all sufficiently fine meshes.
They supply four disjoint strong spin paths, of alternating colors.
The cropping allows their endpoints to be moved by the bounded
lattice distance needed to follow the colored side.

This is an eventual implication on the coupled limit event, so
Fatou's lemma bounds its probability by the lower limit of the
four-strong-arm probabilities at the further relaxed radii. Wu's
Theorem~1.2 gives the interior alternating \emph{spin} exponent
\((16j^2-1)/24\), hence \(21/8\) for four arms. More precisely,
\cite[Equation~(5.4) and Corollary~5.4]{Wu2018} gives the required mesh-first upper
bound \(u^{21/8+o(1)}\) for these annular events, uniformly in a
larger containing square. Combining the two relaxations yields
\eqref{eq:pure-four-passages}; multiplying \(u\) by a fixed constant
does not change its limiting exponent. Square and circular annuli
contain one another after fixed-factor radius changes.

No microscopic inclusion of arbitrary resolved passages in the
strong-arm event was used. Such an inclusion would be false as a
justification at an individual alternating face, where a colored side
of the fixed resolution may have only weak adjacency. Complete
oriented convergence and the extremal side paths are the steps that
remove this convention difference.
\end{proof}

\subsection{Buffered crossings, circuits, and one passage}

\begin{proposition}[Pure buffered RSW consequences]
\label{prop:pure-rsw}
Fix a finite collection of polygonal corridors, each a topological
rectangle with designated opposite end gates, and polygonal annular
bands. Prescribe one color for each required crossing or circuit.
Assume each requirement admits a finite rectangle construction with
positive-width unpinned buffers, and that the buffered constructions
for opposite colors are disjoint. Their simultaneous occurrence has
probability bounded below by a positive constant for all sufficiently
fine meshes. The constant may depend on this fixed geometry, and is
uniform over pins beyond the buffers and ambient plane or torus laws
with common lattice charts.

There are \(C<\infty\), \(\rho>0\) such that, in a buffered bulk
annulus and for \(1\le r<R/C\),
\begin{equation}
 \P(\text{a monochromatic spin arm or a resolved contour passage
                 from }B_r\text{ to }B_R^c)
 \le C(r/R)^\rho.
 \label{eq:pure-one-passage}
\end{equation}
For the monochromatic event either color may be specified or both
allowed. The annular construction may be conditioned on earlier
reveals that remain outside its unpinned test collars. Fixed-factor
changes of radii and bounded lattice drawing corrections change only
the constants.
\end{proposition}

\begin{proof}
Theorem~1 of~\cite{DHN2011} gives critical \(q=2\) FK rectangle
crossing bounds uniform over boundary wirings. Its Section~2.1 gives
the domain Markov property, association, planar duality, and
Edwards--Sokal coloring. A finite chain of overlapping rectangle
crossings, including overlap-square crossings, constructs an FK
connection through a prescribed corridor. The analogous finite
construction gives an FK circuit in a band. All aspect ratios and
the number of rectangles are fixed. A dual circuit in a disjoint
outer collar isolates the constructed cluster from the boundary.
Uniform conditional RSW permits this dual event and the inner open
tests to be imposed in sequence. The isolated cluster then receives
either prescribed spin color with probability \(1/2\). For several
opposite-color constructions keep their isolating collars disjoint;
their clusters can be colored independently. Same-color connections
may merge, and spin FKG combines the finitely many same-color tests.
This proves positivity in the reference free or plane charts.

Mixed-sign spin pins are not arbitrary FK wirings: identifying the
two notions would invalidate this argument. Instead,
Proposition~\ref{prop:pure-buffered} transfers the reference lower
bound through each unpinned collar. Reveal separated constructions
successively; the remaining collars stay free and the lower bounds
remain valid. For a single annular circuit one can first transfer the
individual rectangular tests and then use pure-spin FKG. Consequently
the argument does not require the entire annulus to be a square box
avoiding all interior pins.

For the power bound choose \(m\ge c\log(R/r)-C\) annular bands of
a fixed aspect ratio with separated test collars. Each has a
conditional probability at least \(p>0\) of a strong spin circuit
of either specified color. A circuit of the opposite color blocks a
monochromatic arm. A strong spin circuit of either color also blocks
a contour passage, since disagreement edges cannot cross its
equal-spin bonds. Conditional iteration therefore gives
\(C(1-p)^m\le C'(r/R)^\rho\), with a union bound for two possible
arm colors. This is a spin and contour deduction from FK RSW,
coloring, and buffered comparison; no FK arm exponent is being
identified with a spin arm exponent. Pins touching a corridor's
terminals without a buffer are outside the assertion.
\end{proof}

Together, Proposition~\ref{prop:pure-four-passages} and
Proposition~\ref{prop:pure-rsw} provide, respectively, a bulk exponent
strictly greater than two and a positive one-passage exponent.
Section~\ref{sec:stability} will derive the additional estimate for
two passages constrained to a half-plane, rather than importing a
boundary-arm exponent for a different boundary law.
